\section*{Part V: Accounting: The Language of Enterprise Reality}

\subsection*{38. Why the Founder Must Learn Accounting Early}

A founder need not become a chartered accountant, but they must understand what the financial statements measure and why a growing bank balance can coexist with losses. Accounting is not merely tax compliance. It is a systematic way to distinguish assets from expenses, obligations from income, and performance from liquidity.

3 core statements tell complementary stories:
\begin{enumerate}
  \item The \textbf{income statement} describes revenue and expenses over a defined period.
  \item The \textbf{balance sheet} describes assets, liabilities, and equity at a specific point in time.
  \item The \textbf{cash flow statement} classifies cash movements into operating, investing, and financing activities over a period.
\end{enumerate}
The statement of changes in equity and disclosure notes complete the picture under applicable standards. Official frameworks set out these essential distinctions.\footnote{IFRS Foundation, \emph{Conceptual Framework for Financial Reporting}, revised 2018; U.S. Securities and Exchange Commission, \emph{Beginners' Guide to Financial Statements}.}

\subsection*{39. The Accounting Equation}

At every reporting date, the fundamental identity is
\begin{equation}
\boxed{\text{Assets} = \text{Liabilities} + \text{Equity}}.
\end{equation}
Assets include cash, receivables, certain prepayments, inventory, and qualifying long-lived resources. Liabilities include amounts owed to suppliers, borrowings, and obligations to deliver services after collecting cash. Equity represents the residual interest after liabilities are deducted. An owner funding a company contributes equity or a loan depending on the actual terms; that cash receipt is not customer revenue.

The balance sheet is not a valuation of the founder's dream. An admirable mission, personal expertise, and hoped-for reputation may be strategically valuable without qualifying as accounting assets. Similarly, spending time producing course content does not automatically make it capitalizable intangible property under the applicable framework.

\subsection*{40. The Income Statement: Performance Over a Period}

An income statement normally begins with recognized revenue, subtracts appropriate cost of sales or services, then operating expenses, financing costs, taxes, and other items according to the reporting framework. It answers whether the entity generated accounting income in that reporting period under the applicable recognition and measurement principles.

Revenue is not the same as a signed contract, an invoice, or cash collected. Costs may be incurred before or after payment. Depreciation allocates the cost of qualifying long-lived assets over their useful lives under applicable rules rather than expensing every equipment purchase immediately. A founder who interprets income as a record of bank transfers will misread both profit and cash requirements.

\subsection*{41. Cash-Basis Thinking Versus Accrual Accounting}

Cash-basis tracking records receipts and payments when money moves. It is essential for liquidity management. Accrual accounting recognizes revenue when earned and expenses when incurred under relevant rules, even if cash moves in a different period. Statutory requirements vary by entity and jurisdiction; accounting treatment is a standards question, not a preference for whatever shows more profit.

The key founder insight is practical: \textbf{cash and earned performance are on different clocks}. A business can record a sale today and collect next quarter. It can collect annual subscription cash today while owing 11 more months of access. It can pay for equipment today and recognize depreciation over time. Maintain cash forecasts alongside properly prepared accounts; neither replaces the other.

\subsection*{42. Revenue Recognition: Prepaid Education Is Not Instant Profit}

Under IFRS 15, revenue from a customer contract is recognized to depict the transfer of promised goods or services, applying a 5-step approach:\footnote{IFRS Foundation, \emph{IFRS 15: Revenue from Contracts with Customers}; Government of India, Ministry of Corporate Affairs, \emph{Indian Accounting Standard (Ind AS) 115: Revenue from Contracts with Customers}.}
\begin{enumerate}
  \item Identify the contract with a customer.
  \item Identify the performance obligations in the contract.
  \item Determine the transaction price.
  \item Allocate the transaction price to the performance obligations.
  \item Recognize revenue as or when the entity satisfies each performance obligation.
\end{enumerate}

If a learner prepays for 3 months of instruction, the company may initially owe a service obligation. The cash arrives, but it has not yet earned all the tuition. A \emph{contract liability} arises for advance consideration. Revenue is recognized as the agreed teaching service is provided under the correct contract analysis, not merely when the founder receives a payment notification.

Refund rights, assessment and certification components, multiple distinct services, variable consideration, and cancellations can complicate the schedule. Do not assume every course earns precisely 1-third of the fee each month; use that as a simplified illustration only when a uniform service pattern and contract facts support it.
